sections/02_related_work.tex
\[ %%%%%%%%%%%%%%%%%%%%%%%%%%%%%%%%%%%%%%%%%%%%%%%%%%%%%%%%%%%%%%%%%%%%%%%%%%%%%%%% \section{Related Work} \label{sec:related_work} %%%%%%%%%%%%%%%%%%%%%%%%%%%%%%%%%%%%%%%%%%%%%%%%%%%%%%%%%%%%%%%%%%%%%%%%%%%%%%%% Graph-based learning has emerged as one of the most promising approaches for modeling complex physical systems with irregular spatial structures, heterogeneous interactions, and non-Euclidean geometries. Unlike conventional convolutional neural networks, which are primarily designed for regular grid representations, graph neural networks (GNNs) provide a flexible mathematical framework for representing physical domains as collections of interacting entities connected through relational structures. In scientific computing, this capability is particularly important because many physical systems are naturally described through discretized representations, including finite-element meshes, finite-volume cells, molecular graphs, particle systems, and adaptive computational domains. In these settings, graph representations provide a direct correspondence between computational nodes and physical degrees of freedom, while graph edges encode local interactions governed by physical laws. %%%%%%%%%%%%%%%%%%%%%%%%%%%%%%%%%%%%%%%%%%%%%%%%%%%%%%%%%%%%%%%%%%%%%%%%%%%%%%%% \subsection{Graph Neural Networks for Physical Simulation} \label{sec:gnn_physics} %%%%%%%%%%%%%%%%%%%%%%%%%%%%%%%%%%%%%%%%%%%%%%%%%%%%%%%%%%%%%%%%%%%%%%%%%%%%%%%% The application of graph neural networks to physical simulation has been motivated by the observation that many dynamical systems can be interpreted as interacting collections of local components. A physical state can therefore be represented as an attributed graph \begin{equation} G_t=(V,E,X_t), \end{equation} where $V$ denotes the set of nodes representing physical entities, $E$ represents interaction relationships, and $X_t$ contains time-dependent physical variables associated with each node. The temporal evolution of such a system can be learned through iterative message passing operations that approximate the underlying physical dynamics. One of the earliest influential directions in this field demonstrated that graph neural networks can learn complex physical interactions by treating simulation components as particles connected through relational graphs. These approaches showed that learned message-passing mechanisms can approximate short-term evolution of systems involving fluids, solids, and deformable objects while maintaining generalization across different initial conditions. Subsequent developments introduced graph-based simulators capable of learning continuous dynamics from numerical simulation data. These models typically consist of three fundamental components: an encoder that maps raw physical states into latent representations, a processor that performs repeated graph message passing operations, and a decoder that reconstructs physical variables in the original representation space. Such architectures have achieved impressive results for various computational physics problems, including fluid dynamics, elasticity, and particle-based simulations. The success of graph-based physical simulators can be attributed to several fundamental advantages. First, graph representations naturally preserve locality, which is a central property of many physical interactions. Second, message-passing mechanisms provide permutation invariance, allowing models to operate independently of arbitrary node ordering. Third, graph architectures can accommodate irregular geometries and adaptive discretizations that are difficult to handle using traditional convolutional neural networks. However, despite these advantages, conventional graph neural simulators remain primarily data-driven. The physical laws governing the simulated system are usually represented implicitly through training examples rather than explicitly incorporated into the computational mechanism. As a consequence, the learned dynamics may reproduce observed trajectories while failing to preserve fundamental physical properties outside the training distribution or over extended forecasting horizons. A critical limitation of existing GNN-based simulators is error accumulation during autoregressive prediction. Since future states are generated recursively using previous model outputs, small approximation errors can propagate through multiple message-passing iterations. In nonlinear systems, these errors may amplify rapidly, leading to unstable trajectories and violation of physical constraints. This issue is especially pronounced in chaotic systems, turbulent flows, and multi-scale physical processes where long-term evolution is highly sensitive to initial perturbations. Several approaches have attempted to improve the physical reliability of graph neural simulators by introducing conservation constraints, energy-based regularization, or hybrid numerical-learning strategies. While these methods represent important progress, physical knowledge is often incorporated as an additional training objective rather than as an intrinsic component of the computational architecture. Furthermore, many existing graph-based models focus primarily on short-term prediction accuracy and provide limited mechanisms for monitoring stability, uncertainty propagation, or physical admissibility during long-horizon inference. This limitation restricts their application in scenarios requiring reliable forecasting, such as digital twins, autonomous control systems, and real-time scientific decision-making. These observations motivate the development of physics-guided graph architectures in which physical principles influence not only optimization objectives but also the internal mechanisms of representation learning, information propagation, and temporal state evolution. The PHYSGEN framework builds upon the strengths of graph-based scientific learning while addressing these limitations through a multi-level physics guidance strategy designed specifically for stable long-horizon dynamical system prediction. \] 

